\section{2026 展望：训练的不只是模型，而是整个系统}

\subsection{从模型中心走向系统中心}
panel 后段有一句非常值得记住的话：未来不只是训练模型，而是在训练一个系统。这里的系统包括模型前后的渲染、推理、工作流、环境构造、数据回流和部署接口。它意味着 ``模型'' 不再是唯一优化对象，而只是整条链路里最核心、但不是唯一的一环。

\begin{importantbox}{2026 年高概率继续扩张的方向}
\begin{itemize}
  \item 推理与训练的结合会继续加深，训推分离会从架构问题变成 workload 问题；
  \item Agent environment 的构造会成为重要工程能力；
  \item 多模态、RL 与 Agent workload 会继续推动系统异构化；
  \item 历史上的分布式系统、数据库与容错经验，会被重新引入大模型 Infra。
\end{itemize}
\end{importantbox}

\subsection{Infra 的终局不是隐藏，而是成为默认能力}
从 panel 的语气能听出来，嘉宾们并不追求让用户直接感知 Infra，而是希望它最终退到后台，成为系统默认具备的支撑能力。也就是说，当我们讨论 ``Agent 产品好不好用'' 时，很多真实决定因素其实都藏在 rollout 引擎、缓存策略、调度系统、环境接口和故障恢复里。

\begin{table}[H]
\centering
\begin{tabular}{p{2.8cm}p{4.2cm}p{4.4cm}}
\toprule
\textbf{方向} & \textbf{系统目标} & \textbf{Infra 侧关键问题} \\
\midrule
推理继续集群化 & 更高吞吐、更低尾延迟 & PD 分离、专家并行、跨节点调度 \\
RL / Agent 规模化 & 更长 rollout、更大环境、更复杂反馈 & 训推结合、environment 接口、reward 回流 \\
异构训练系统 & 支持更多 workload 与硬件组合 & 模块化抽象、故障恢复、精度管理 \\
产业落地 & 让模型真正接入应用链路 & 文件系统、权限、workflow、稳定性 \\
\bottomrule
\end{tabular}
\caption{把整场讨论压缩成 ``Infra 下一步要接住什么''}
\end{table}

\subsection{本章小结}
这场 Infra 专题最终落到一个非常明确的判断：大模型系统的下一个阶段，不是单纯堆更多卡，而是让推理、训练、环境和工作流形成一个真正可扩展的整体系统。这也是 Infra 在 Agent 时代重新变成主线的根本原因。

